% GENERATED FILE. Edit canonical lessons, then run npm run generate:books.
% canonical-source-hash: fnv1a64:efaaa64a973c9248
% canonical-lessons: ES-C693-olivo, ES-C693-pino, ES-C693-arbusto, ES-C693-vegetacion, ES-C693-roble

\chapter{Árboles y plantas}
\label{ch:es-arboles-y-plantas}
\hlchaptermodality{693}

\begin{chapteropening}
\textbf{By the end of this chapter:} \emph{I can name trees and plants on a walk: olive trees, pines, oaks, bushes and plant life.}

Complete the last lesson of chapter 693: I can name trees and plants on a walk: olive trees, pines, oaks, bushes and plant life.
\end{chapteropening}

\section[el olivo]{el olivo — the olive tree}

\label{lesson:ES-C693-olivo}



\begin{quote}
Before the new one: say the Spanish for the well, then the Spanish for to flower.
\end{quote}

\subsection*{You'll want to know: el olivo}
\textbf{el olivo} — \textquotedblleft{}the olive tree\textquotedblright{}.

\begin{quote}
En el campo hay miles de olivos.
\end{quote}

\begin{quote}
\emph{There are thousands of olive trees in the countryside.}
\end{quote}

The fruit is \emph{la aceituna}; the oil is \emph{el aceite de oliva}.

\begin{cousinweb}
\emph{olivo} is inherited from Vulgar Latin \emph{olīvus}, from Latin \emph{olīva}, which comes from Greek \emph{elaía} \textquotedblleft{}olive\textquotedblright{} (EtymDB).
\end{cousinweb}

\begin{grammarlens}[title={What you've built}]
One more word for telling, explaining and describing at B1.
\end{grammarlens}

\subsection*{Your turn}
\begin{itemize}
\raggedright
  \item \emph{Say it:} \emph{el olivo}
  \item \emph{Say it:} \emph{el olivo}, once more
  \item \emph{Say it:} a sentence of your own with \emph{el olivo}
  \item \emph{Recall:} say the Spanish for the well, then the Spanish for to flower, then say \emph{el olivo} again
\end{itemize}

\begin{tcolorbox}[breakable,colback=teal!4,colframe=teal!35!black,title={Before you move on}]
What does el olivo mean? (\textquotedblleft{}The olive tree\textquotedblright{}.) Say it once more.
\end{tcolorbox}
\section[el pino]{el pino — the pine tree}

\label{lesson:ES-C693-pino}



\begin{quote}
Before the new one: say the Spanish for to flower, then the Spanish for the olive tree.
\end{quote}

\subsection*{You'll want to know: el pino}
\textbf{el pino} — \textquotedblleft{}the pine tree\textquotedblright{}.

\begin{quote}
El camino pasa por un bosque de pinos.
\end{quote}

\begin{quote}
\emph{The path goes through a pine forest.}
\end{quote}

\emph{La piña} is the pine cone, and also a pineapple.

\begin{cousinweb}
\emph{pino} is inherited from Latin \emph{pīnus} \textquotedblleft{}a pine tree\textquotedblright{} (EtymDB).
\end{cousinweb}

\begin{grammarlens}[title={What you've built}]
One more word for telling, explaining and describing at B1.
\end{grammarlens}

\subsection*{Your turn}
\begin{itemize}
\raggedright
  \item \emph{Say it:} \emph{el pino}
  \item \emph{Say it:} \emph{el pino}, once more
  \item \emph{Say it:} a sentence of your own with \emph{el pino}
  \item \emph{Recall:} say the Spanish for to flower, then the Spanish for the olive tree, then say \emph{el pino} again
\end{itemize}

\begin{tcolorbox}[breakable,colback=teal!4,colframe=teal!35!black,title={Before you move on}]
What does el pino mean? (\textquotedblleft{}The pine tree\textquotedblright{}.) Say it once more.
\end{tcolorbox}
\section[el arbusto]{el arbusto — the bush, the shrub}

\label{lesson:ES-C693-arbusto}



\begin{quote}
Before the new one: say the Spanish for the olive tree, then the Spanish for the pine tree.
\end{quote}

\subsection*{You'll want to know: el arbusto}
\textbf{el arbusto} — \textquotedblleft{}the bush, the shrub\textquotedblright{}.

\begin{quote}
Hay arbustos a los dos lados del camino.
\end{quote}

\begin{quote}
\emph{There are bushes on both sides of the path.}
\end{quote}

Smaller than \emph{un árbol}.

\begin{cousinweb}
\emph{arbusto} is a learned borrowing from Latin \emph{arbustum} \textquotedblleft{}a place planted with trees\textquotedblright{} (Lewis \& Short; EtymDB), from \emph{arbor} \textquotedblleft{}tree\textquotedblright{}.
\end{cousinweb}

\begin{grammarlens}[title={What you've built}]
One more word for telling, explaining and describing at B1.
\end{grammarlens}

\subsection*{Your turn}
\begin{itemize}
\raggedright
  \item \emph{Say it:} \emph{el arbusto}
  \item \emph{Say it:} \emph{el arbusto}, once more
  \item \emph{Say it:} a sentence of your own with \emph{el arbusto}
  \item \emph{Recall:} say the Spanish for the olive tree, then the Spanish for the pine tree, then say \emph{el arbusto} again
\end{itemize}

\begin{tcolorbox}[breakable,colback=teal!4,colframe=teal!35!black,title={Before you move on}]
What does el arbusto mean? (\textquotedblleft{}The bush, the shrub\textquotedblright{}.) Say it once more.
\end{tcolorbox}
\section[la vegetación]{la vegetación — the vegetation, the plant life}

\label{lesson:ES-C693-vegetacion}



\begin{quote}
Before the new one: say the Spanish for the pine tree, then the Spanish for the bush.
\end{quote}

\subsection*{You'll want to know: la vegetación}
\textbf{la vegetación} — \textquotedblleft{}the vegetation, the plant life\textquotedblright{}.

\begin{quote}
La vegetación es muy verde en el norte.
\end{quote}

\begin{quote}
\emph{The vegetation is very green in the north.}
\end{quote}

A general word for all the plants of a place.

\begin{cousinweb}
\emph{vegetación} is a learned borrowing from Late Latin \emph{vegetātiō} \textquotedblleft{}an enlivening, a quickening\textquotedblright{} (Lewis \& Short), from \emph{vegetāre} \textquotedblleft{}to enliven\textquotedblright{}.
\end{cousinweb}

\begin{grammarlens}[title={What you've built}]
One more word for telling, explaining and describing at B1.
\end{grammarlens}

\subsection*{Your turn}
\begin{itemize}
\raggedright
  \item \emph{Say it:} \emph{la vegetación}
  \item \emph{Say it:} \emph{la vegetación}, once more
  \item \emph{Say it:} a sentence of your own with \emph{la vegetación}
  \item \emph{Recall:} say the Spanish for the pine tree, then the Spanish for the bush, then say \emph{la vegetación} again
\end{itemize}

\begin{tcolorbox}[breakable,colback=teal!4,colframe=teal!35!black,title={Before you move on}]
What does la vegetación mean? (\textquotedblleft{}The vegetation, the plant life\textquotedblright{}.) Say it once more.
\end{tcolorbox}
\section[el roble]{el roble — the oak tree}

\label{lesson:ES-C693-roble}



\begin{quote}
Before the new one: say the Spanish for the bush, then the Spanish for the vegetation.
\end{quote}

\subsection*{You'll want to know: el roble}
\textbf{el roble} — \textquotedblleft{}the oak tree\textquotedblright{}.

\begin{quote}
Comimos a la sombra de un roble.
\end{quote}

\begin{quote}
\emph{We ate in the shade of an oak tree.}
\end{quote}

\emph{Fuerte como un roble} — as strong as an ox.

\begin{cousinweb}
\emph{roble} is inherited from Latin \emph{rōbur} \textquotedblleft{}an oak tree; hardness, strength\textquotedblright{} (EtymDB).
\end{cousinweb}

\begin{grammarlens}[title={What you've built}]
One more word for telling, explaining and describing at B1.
\end{grammarlens}

\subsection*{Your turn}
\begin{itemize}
\raggedright
  \item \emph{Say it:} \emph{el roble}
  \item \emph{Say it:} \emph{el roble}, once more
  \item \emph{Say it:} a sentence of your own with \emph{el roble}
  \item \emph{Recall:} say the Spanish for the bush, then the Spanish for the vegetation, then say \emph{el roble} again
\end{itemize}

\begin{tcolorbox}[breakable,colback=teal!4,colframe=teal!35!black,title={Before you move on}]
What does el roble mean? (\textquotedblleft{}The oak tree\textquotedblright{}.) Say it once more.
\end{tcolorbox}
